\begin{definition}[Geometric synchronized scheme]
  \label{def:geometric-synchronized-scheme}
  A geometric synchronized scheme for \(U\) consists of
  \(a\in A^{J_{\mathrm r}}\), an integer \(k\ge1\), integers
  \(c\ge1,b\ge2\), and words
  \[
    u_1,\ldots,u_{k+1},\quad
    v_1,\ldots,v_k,\quad w_1,\ldots,w_k,\quad
    s_1,\ldots,s_k\in A^\ast,
  \]
  such that
  \[
    W(t)=a u_1v_1^tw_1s_1^tu_2\cdots
      u_kv_k^tw_ks_k^tu_{k+1}\in\mathcal R_U
      \qquad(t\ge0),
  \]
  \[
    \sum_{j=1}^k(\abs{v_j}+\abs{s_j})>0,
    \qquad \val_U(W(t))=cb^t\quad(t\ge0).
  \]
\end{definition}

\begin{lemma}[Arithmetic input of Evertse]
  \label{lem:evertse-quotient-support-input}
  Let \(K\) be a number field. For \(x\in K^\times\), let \(P_K(x)\)
  be the largest norm of a finite prime ideal \(\mathfrak p\) for which
  \(\operatorname{ord}_{\mathfrak p}(x)\ne0\), with \(P_K(x)=1\) when
  \(x\) is a unit. If \(Z=(Z(n))_{n\ge0}\) is a nondegenerate linear
  recurrence over \(K\) whose minimal shift-annihilating polynomial has
  at least two distinct nonzero roots, then
  \[
    \lim_{\substack{r\to\infty\\ r>s\ge0,\ Z(s)\ne0}}
      P_K\!\left(\frac{Z(r)}{Z(s)}\right)=\infty.
  \]
  In particular, a rational-integer recurrence with no zero terms, all
  supported on one finite set of rational primes, cannot be both
  nondegenerate and have at least two characteristic roots.
\end{lemma}

\begin{proof}
  The displayed uniform limit is exactly
  \cite[Theorem~3]{Evertse1984}. For the consequence, let \(S\) be the
  finite rational-prime support and choose a number field \(K\) containing
  the characteristic roots and coefficients. If \(Z(r),Z(s)\ne0\), the
  principal fractional ideal generated by \(Z(r)/Z(s)\) is supported
  only at prime ideals above primes in \(S\). Every prime ideal
  \(\mathfrak p\) above \(p\in S\) has
  \[
    \operatorname N\mathfrak p=p^{f(\mathfrak p/p)}
      \le p^{[K:\mathbb Q]}.
  \]
  Consequently
  \[
    P_K\!\left(\frac{Z(r)}{Z(s)}\right)
      \le \max\bigl(\{p^{[K:\mathbb Q]}:p\in S\}\cup\{1\}\bigr)
  \]
  for every \(r,s\) with nonzero denominator. This fixed bound
  contradicts Evertse's limit when there are at least two characteristic
  roots. Here and below, ``characteristic roots of a sequence'' always
  means the roots of its minimal shift-annihilating polynomial, not roots
  introduced by a nonminimal recurrence witness.
\end{proof}

\begin{lemma}[Injective \(S\)-unit recurrences]
  \label{lem:s-unit-increasing-recurrence-rigidity}
  Let \(Y=(Y(n))_{n\ge0}\) be a pairwise distinct sequence of positive
  rational integers satisfying a linear recurrence. If all prime
  divisors of all \(Y(n)\) belong to one finite set, then there are
  integers \(r\ge0,h\ge1,c\ge1\), and \(b\ge2\) such that
  \[
    Y(r+hn)=cb^n\qquad(n\ge0).
  \]
\end{lemma}

\begin{proof}
  Write the standard exponential-polynomial expansion, after discarding
  the contribution of the zero characteristic root and a finite initial
  segment, as
  \[
    Y(n)=\sum_{i=1}^m f_i(n)\alpha_i^n,
  \]
  where the \(\alpha_i\) are distinct nonzero algebraic numbers and the
  nonzero \(f_i\) are algebraic polynomials. Partition the roots by
  \(\alpha_i\sim\alpha_j\) when \(\alpha_i/\alpha_j\) is a root of
  unity. Choose \(h_0\ge1\) divisible by the orders of all such
  quotients. For each \(j\in\{0,\ldots,h_0-1\}\), substitution of
  \(n=j+h_0t\) and collection within the equivalence classes gives
  \[
    Y(j+h_0t)=\sum_C g_{C,j}(t)\beta_C^t,
    \qquad \beta_C=\alpha_i^{h_0}\quad(i\in C).
  \]
  Delete the identically zero summands. Distinct remaining \(\beta_C\)
  have quotient that is not a root of unity. Thus every residue
  subsequence is a nonzero nondegenerate recurrence; it is nonzero
  because all of its terms are positive.

  Every term of every residue subsequence is a nonzero rational integer
  supported on the prescribed finite set. If a residue subsequence had
  at least two characteristic roots, the consequence in
  \Cref{lem:evertse-quotient-support-input} would contradict the fixed
  support of all its term quotients. Hence every residue subsequence has
  exactly one characteristic root and therefore has the form
  \[
    Z(n)=f(n)\lambda^n
  \]
  with one characteristic root.

  Choose one residue subsequence and shift its index if necessary so
  that the displayed expansion holds from \(n=0\), and put
  \(e=\deg f\). Let \(E\) denote the forward shift. Over
  \(\overline{\mathbb Q}\), finite differences give
  \[
    (E-\lambda)^{e+1}Z=0,
    \qquad (E-\lambda)^eZ\ne0,
  \]
  because the \(e\)-th difference of the nonzero degree-\(e\) polynomial
  \(f\) is a nonzero constant. Hence the monic minimal shift-annihilator
  over \(\overline{\mathbb Q}\) is exactly
  \((X-\lambda)^{e+1}\), and the Hankel rank of \(Z\) over that field is
  \(e+1\).

  The Hankel matrix has rational entries, and a rational matrix has the
  same rank over \(\mathbb Q\) and over \(\overline{\mathbb Q}\).
  Therefore a minimal dependence among its shift columns gives a monic
  polynomial \(P\in\mathbb Q[X]\) of the same degree \(e+1\). After
  extending scalars, minimality and equality of degrees force
  \[
    P(X)=(X-\lambda)^{e+1}.
  \]
  For every
  \(\sigma\in\operatorname{Gal}(\overline{\mathbb Q}/\mathbb Q)\), the
  polynomial \(P\) is fixed and has only the root \(\lambda\); thus
  \(\sigma(\lambda)=\lambda\). Hence \(\lambda\in\mathbb Q\). Now
  \(f(n)=Z(n)/\lambda^n\in\mathbb Q\)
  for every \(n\), and interpolation from \(\deg(f)+1\) integer
  arguments gives \(f\in\mathbb Q[X]\). If \(f\) were nonconstant, clear its
  denominators to obtain a nonconstant \(g\in\mathbb Z[X]\). The
  identity \(Z(n)=f(n)\lambda^n\) would confine every prime divisor of
  every nonzero \(g(n)\) to a fixed finite set. Indeed, if
  \(\lambda=a/b\) is in lowest terms and \(g=Df\), then
  \[
    g(n)a^n=Db^nZ(n).
  \]
  Thus that fixed set is the union of the prescribed prime set and the
  prime divisors of \(Dab\). This contradicts Schur's theorem that a
  nonconstant integer polynomial has values divisible by infinitely
  many primes; see, for example, \cite{HardyWright2008}. Hence \(f\) is
  constant.

  We have therefore \(Z(n)=c\lambda^n\), where
  \(c=Z(0)\) is a positive integer. If \(\lambda=a/b\) in lowest terms,
  integrality for every \(n\) forces \(b^n\mid c\) for every \(n\), so
  \(b=1\). Positivity gives \(\lambda>0\), and pairwise distinctness
  excludes \(\lambda=1\). Hence \(\lambda\in\mathbb Z_{\ge2}\).
  Translating the chosen residue and
  index shift back to \(Y\) supplies the asserted \(r\) and \(h\).
\end{proof}

\begin{theorem}[Geometric-ray characterization]
  \label{thm:bounded-outside-support-geometric-ray}
  For a canonical eventually recurrent numeration system satisfying
  \textup{(U1)}, \textup{(U2)}, and \textup{(U3)}, the following are
  equivalent.
  \begin{enumerate}[label=\textup{(\roman*)}]
    \item There are a finite \(k\) and an infinite \(k\)-MCFL
    \(L\subseteq\mathcal R_U\) such that
    \[
      \sup_{w\in L}\omega_{S_{\mathrm r}^c}(\val_U(w))<\infty.
    \]
    \item The system admits a geometric synchronized scheme in the
    sense of \Cref{def:geometric-synchronized-scheme}.
    \item The language \(\mathcal R_U\) contains an infinite linear MCFL
    \(L_0=\{W(t):t\ge0\}\) of finite fan-out for which there are
    \(c\ge1,b\ge2\) satisfying
    \[
      \val_U(W(t))=cb^t\qquad(t\ge0).
    \]
  \end{enumerate}
\end{theorem}

\begin{proof}
  Suppose first that \textup{(i)} holds. Use the minimal tail recurrence
  as the witness in
  \Cref{thm:mcfl-synchronized-congruence-orbit}, and write its resulting
  ray and pairwise distinct value sequence as \(W(t)\) and \(N(t)\).
  Let
  \[
    K_0=\max_{t\ge0}\omega_{S_{\mathrm r}^c}(N(t)),
  \]
  and choose \(t_0\) attaining the maximum. Let \(T\) be the set of
  primes outside \(S_{\mathrm r}\) that divide \(N(t_0)\), and put
  \(q=\prod_{\ell\in T}\ell\). If \(T\ne\varnothing\), take the return
  time \(H=H(q)\) from
  \Cref{thm:mcfl-synchronized-congruence-orbit}; if \(T=\varnothing\),
  put \(H=1\). Every prime in \(T\) divides every \(N(t_0+nH)\). A prime
  divisor outside \(S_{\mathrm r}\cup T\) would therefore give at least
  \(K_0+1\) distinct prime divisors outside \(S_{\mathrm r}\), contrary
  to the definition of \(K_0\). Thus all terms of
  \[
    Y(n)=N(t_0+nH)
  \]
  are supported on the fixed finite set \(S_{\mathrm r}\cup T\).

  To make the termwise-product closure explicit, put
  \(n_0=d_{\mathrm r}+1\), list the \(2k\) powered affine blocks as
  \(B_1,\ldots,B_{2k}\), and form
  \[
    K=B_1\otimes\cdots\otimes B_{2k}.
  \]
  Since \(K^t=B_1^t\otimes\cdots\otimes B_{2k}^t\), expansion of the
  interspersed constant block matrices expresses \(N(t)\) as a fixed
  linear functional of \(K^t\). Cayley--Hamilton therefore gives a
  recurrence of order at most
  \[
    R=n_0^{2k}=(d_{\mathrm r}+1)^{2k}.
  \]
  Arithmetic subsequences of a recurrence are recurrences, so \(Y\) is a
  pairwise distinct positive integer recurrence. By
  \Cref{lem:s-unit-increasing-recurrence-rigidity}, there are
  \(r\ge0,h\ge1,c\ge1,b\ge2\) such that
  \(N(t_\ast+h_\ast n)=cb^n\), where
  \(t_\ast=t_0+Hr\) and \(h_\ast=Hh\).

  In the synchronized factorization, replace
  \[
    u_j\ \text{by}\ u_jv_j^{t_\ast},\qquad
    v_j\ \text{by}\ v_j^{h_\ast},\qquad
    w_j\ \text{by}\ w_js_j^{t_\ast},\qquad
    s_j\ \text{by}\ s_j^{h_\ast}.
  \]
  This absorbs the fixed powers and parametrizes exactly the subray
  \(W(t_\ast+h_\ast n)\). Its total pumped length remains positive, so
  it is a geometric synchronized scheme. This proves
  \textup{(i)}\(\Rightarrow\)\textup{(ii)}.

  The implication \textup{(ii)}\(\Rightarrow\)\textup{(iii)} is
  \Cref{lem:synchronized-ray-linear-mcfg}. Finally, under
  \textup{(iii)}, the language \(L_0\) is an infinite MCFL and
  \[
    \omega_{S_{\mathrm r}^c}(cb^t)\le\omega(cb)
    \qquad(t\ge0),
  \]
  proving \textup{(i)}.
\end{proof}

\begin{lemma}[Prime support of the geometric ratio]
  \label{lem:geometric-ratio-tail-support}
  In every geometric synchronized scheme
  \(\val_U(W(t))=cb^t\), each prime divisor of \(b\) divides the
  constant coefficient \(c_0\) of the minimal tail recurrence.  Equivalently,
  \[
    \operatorname{rad}(b)\mid\operatorname{rad}(\abs{c_0}).
  \]
\end{lemma}

\begin{proof}
  Suppose that a prime \(p\mid b\) does not divide \(c_0\), put
  \(e=v_p(c)\), and take the allowed modulus \(q=p^{e+1}\).  Every
  nonempty block matrix among the \(M_{v_j}\) and \(M_{s_j}\) is invertible
  modulo \(q\) by \Cref{lem:affine-recurrence-block-action}.  A common
  multiple \(H\ge1\) of their orders gives, by the matrix product defining
  the scheme exactly as in the proof of
  \Cref{thm:mcfl-synchronized-congruence-orbit},
  \[
    \val_U(W(t+H))\equiv\val_U(W(t))\pmod q
    \qquad(t\ge0).
  \]
  At \(t=0\) this says
  \[
    cb^H\equiv c\pmod {p^{e+1}}.
  \]
  But \(p\mid b\) implies \(p\nmid b^H-1\), so
  \(v_p(c(b^H-1))=e\), contradicting the displayed congruence.
\end{proof}

We call the following finite data an \emph{effective presentation}: the
digit alphabet, initial values and recurrence coefficients with a valid
threshold, a DFA for \(\mathcal R_U\), and the minimal tail recurrence with
a valid threshold. The presentation is assumed to satisfy
\textup{(U1)}--\textup{(U3)}, including uniqueness of canonical
representations; recognizing valid presentations is not part of the
algorithm.

\begin{corollary}[Semidecidability]
  \label{cor:bounded-outside-support-semidecidable}
  From effective presentations, the set of positive instances of the property in
  \Cref{thm:bounded-outside-support-geometric-ray}\textup{(i)} is
  recursively enumerable, equivalently it has a \(\Sigma^0_1\) upper
  bound.
\end{corollary}

\begin{proof}
  Enumerate \(k\ge1\), \(a\in A^{J_{\mathrm r}}\), and all finite tuples
  of words occurring in
  \Cref{def:geometric-synchronized-scheme}, rejecting tuples of zero
  total pumped length. For a fixed tuple, whether the DFA accepts
  \(W(t)\) for every \(t\ge0\) is decidable: the tuple of powers of the
  finitely many pumped-block transformations is eventually periodic in
  the finite transition monoid, and its preperiod and period are found
  by exhaustive iteration.

  The affine matrices construct an exact linear recurrence for
  \[
    N(t)=\val_U(W(t)).
  \]
  With \(d_{\mathrm r}\) the minimal tail order, the Kronecker construction
  in the proof of
  \Cref{thm:bounded-outside-support-geometric-ray} has dimension
  \[
    R=(d_{\mathrm r}+1)^{2k};
  \]
  it is computable from the candidate and is an upper bound for the
  recurrence order of \(N(t)\). For a tuple accepted at every parameter,
  compute \(c=N(0)\). The only
  possible integer ratio is \(b=N(1)/N(0)\). Check that
  \(b\in\mathbb Z_{\ge2}\), and then decide whether the recurrence
  \(N(t+1)-bN(t)\) is identically zero by checking its first \(R\) terms;
  Cayley--Hamilton makes those \(R\) zero values sufficient. Thus each
  candidate scheme is decidable, and the enumeration halts exactly on
  positive instances by
  \Cref{thm:bounded-outside-support-geometric-ray}.
\end{proof}

\begin{remark}[Decision boundary]
  \label{rem:general-decision-interface}
  A decision procedure would follow from a computable bound \(B\), derived
  from the input DFA and minimal tail recurrence, such that every positive
  instance has a geometric synchronized scheme whose fan-out and total
  block length are at most \(B\). No such bound is known. Replacing a long
  DFA loop by a shorter word with the same endpoints can change its affine
  recurrence matrix and destroy the identity
  \(\val_U(W(t+1))=b\val_U(W(t))\). The positive semidecision therefore
  establishes neither decidability nor undecidability of the general
  problem.
\end{remark}

\begin{theorem}[Arbitrarily deep quotient congruences]
  \label{thm:unit-mcfl-deep-congruence-chain}
  Suppose \(\abs{a_0}=1\), and let \(L\subseteq\mathcal R_U\) be an
  infinite MCFL. For arbitrary sequences \(M_i\ge1\) and \(r_i\ge1\),
  \(L\) contains a strictly increasing sequence \(N_i\) such that
  \[
    N_{i+1}=N_iQ_i,\qquad Q_i>1,\qquad
    Q_i\equiv1\pmod {M_iN_i^{r_i}}.
  \]
  If \(M_i\) is divisible by every prime at most \(B_i\), then every
  prime divisor of \(Q_i\) exceeds \(B_i\).
\end{theorem}

\begin{proof}
  Take the fixed orbit \(N(t)\) from
  \Cref{thm:mcfl-synchronized-congruence-orbit}, discarding finitely many
  initial parameters so that \(N(t)\ge2\). Having selected
  \(N_i=N(t_i)\), apply the synchronized congruence with
  \[
    q_i=M_iN_i^{r_i+1}.
  \]
  The values \(N(t_i+mH(q_i))\), \(m\ge1\), are pairwise distinct and all
  are congruent to \(N_i\) modulo \(q_i\). Only finitely many positive
  integers are at most \(N_i\), so choose \(m_i\ge1\) for which
  \(N(t_i+m_iH(q_i))>N_i\), and set
  \(t_{i+1}=t_i+m_iH(q_i)\). The congruence gives
  \[
    N_{i+1}=N_i+c_iM_iN_i^{r_i+1}
           =N_i(1+c_iM_iN_i^{r_i})
  \]
  for an integer \(c_i\ge1\). This proves the first assertion. Under
  the last hypothesis, \(Q_i\equiv1\pmod p\) for every prime
  \(p\le B_i\), so none of those primes divides \(Q_i\).
\end{proof}

\begin{corollary}[Induced full divisibility tree]
  \label{thm:unit-mcfl-divisibility-tree}
  Under the hypotheses of
  \Cref{thm:unit-mcfl-deep-congruence-chain}, there is an injection
  \[
    \Phi:\NN^{<\omega}\longrightarrow L
  \]
  such that
  \[
    \val_U(\Phi(\sigma))\mid\val_U(\Phi(\tau))
      \quad\Longleftrightarrow\quad \sigma\preceq\tau,
  \]
  where \(\preceq\) denotes the prefix order. The quotients attached to
  distinct edges can be chosen pairwise coprime. Moreover, given an
  arbitrary prescribed threshold for each edge, every prime divisor of
  its quotient can be required to exceed that threshold.
\end{corollary}

\begin{proof}
  Use the fixed pairwise-distinct orbit from
  \Cref{thm:mcfl-synchronized-congruence-orbit}. Here is an explicit
  parent-before-child enumeration. If \(p_i\) is the \(i\)-th prime, encode
  \(\sigma=(\sigma_0,\ldots,\sigma_{m-1})\) by
  \[
    \kappa(\sigma)=2^m\prod_{i=0}^{m-1}p_{i+2}^{\sigma_i+1}.
  \]
  Unique factorization makes \(\kappa\) injective, and
  \(\kappa(\sigma^\frown j)>\kappa(\sigma)\). Enumerating the nonroot
  nodes, equivalently their incoming edges, by increasing \(\kappa\)
  therefore has order type \(\omega\) and places every parent before each
  child. Map
  the root to any orbit word of value \(R_0\ge2\). Suppose finitely many
  edges have been assigned, and consider the next edge from a parent of
  value \(P\). Let \(R\) be divisible by \(R_0\), by every previously
  assigned edge quotient, and by every prime not exceeding the prescribed
  threshold for that edge. Apply the synchronized orbit congruence
  modulo \(P^2R\) at the parameter of the parent. Positive multiples of
  the return time give pairwise distinct positive values in the indicated
  congruence class. Only finitely many are at most \(P\), so some multiple
  gives an unused orbit word of value \(P'>P\) with
  \[
    P'\equiv P\pmod {P^2R}.
  \]
  Hence \(P'=PQ\), where \(Q>1\) and \(Q\equiv1\pmod {PR}\). In
  particular, \(Q\) is coprime to the root and to every earlier edge
  quotient, and no prime up to the prescribed threshold divides \(Q\).

  Recursion defines \(\Phi\) on every node. The value at a node is the
  root value times the product of the edge quotients along its path.
  Since the root and all distinct edge quotients are pairwise coprime,
  the value at \(\sigma\) divides the value at \(\tau\) exactly when the
  root-to-\(\sigma\) edge set is contained in the root-to-\(\tau\) edge
  set. For two nodes of a rooted tree, containment of their root-path edge
  sets holds exactly when the first node is an ancestor of the second;
  here this is precisely \(\sigma\preceq\tau\). This also proves injectivity.
\end{proof}

\subsection{Consequences for standard greedy systems}

\begin{corollary}
  \label{thm:recurrent-numeration-prime-cf-immunity}
  Let \((U,A,\mathcal R_U)\) satisfy \textup{(U1)}, \textup{(U2)}, and
  \textup{(U3)}, and put
  \[
    \mathcal P_U
      :=\{w\in\mathcal R_U:\val_U(w)\text{ is prime}\}.
  \]
  Then \(\mathcal P_U\) is MCF-immune, and therefore CF-immune,
  REG-immune, and not recognizable by a finite automaton over \(A\).

  If, in addition, \(\abs{a_0}=1\), then every infinite multiple
  context-free
  \(L\subseteq\mathcal R_U\) has unbounded
  \(\omega(\val_U(w))\). More precisely, \(L\) contains a
  coprime-quotient chain \(N_{i+1}=N_iQ_i\) with
  \(Q_i>1\) and \(Q_i\equiv1\pmod {N_i}\). Thus the bounded-\(\omega\),
  bounded-\(\Omega\), and finite-prime-support languages of a unit
  system are MCF-immune whenever they are infinite. Without the unit
  hypothesis, every bounded-\(\Omega\) language is still MCF-immune.
\end{corollary}

\begin{proof}
  The MCF-immunity assertions follow from
  \Cref{thm:nonunit-mcfl-escape-dichotomy}, including its
  bounded-\(\Omega\) conclusion. In a unit system \(S\) is empty, so
  that theorem gives unbounded \(\omega\). Taking
  \(M_i=r_i=1\) in
  \Cref{thm:unit-mcfl-deep-congruence-chain} gives the stated chain.
  Regular languages are context-free and context-free languages are
  MCFLs. Finally, \textup{(U2)} and Euclid's theorem make
  \(\mathcal P_U\) infinite.
\end{proof}

\paragraph{Greedy length interval.}
For a standard greedy linear numeration basis with \(U_0=1\), a positive
integer has a canonical greedy word of length \(m\) if and only if it lies
in the half-open interval
\[
  [U_{m-1},U_m).
\]
Indeed, the greedy algorithm chooses the least threshold \(U_m\) strictly
larger than the represented integer; minimality of \(m\) gives the lower
bound, and the choice of \(U_m\) gives the upper bound. This is a property
of the standard greedy subclass, not an assumption on the ambient
canonical eventually recurrent systems.

\begin{corollary}
  \label{cor:linear-perron-prime-cf-immunity}
  Let \(U=(U_n)_{n\ge0}\) be a standard greedy linear numeration basis in
  the following precise sense: \(U_0=1\), the \(U_n\) are strictly
  increasing integers with bounded quotients, \(\mathcal R_U\) is the LSD-first
  reversal of the usual greedy representation language, and \(U\)
  has minimal tail recurrence polynomial equal to the minimal polynomial
  of a Perron number \(\beta\). Thus
  \(\beta>1\) is a real algebraic integer and every other conjugate
  \(\beta'\) satisfies \(\abs{\beta'}<\beta\); compare
  \cite{LindMarcus1995,Rigo2014Vol2}. Then its prime
  representation language is MCF-immune, hence CF-immune, REG-immune and
  nonrecognizable. The same statements hold in the usual MSD-first
  convention.

  If the dominant Perron number is a unit and the chosen recurrence has
  trailing coefficient \(a_0=\pm1\), then the full coprime-quotient and
  bounded-prime-support MCF-immunity conclusions of
  \Cref{thm:recurrent-numeration-prime-cf-immunity} hold.
\end{corollary}

\begin{proof}
  Greedy expansion gives a unique representation of every positive
  integer, and its length-\(m\) words represent values between successive
  place-value thresholds; thus \textup{(U1)} and \textup{(U2)} hold. The defining
  Perron recurrence gives \textup{(U3)}. Apply
  \Cref{thm:recurrent-numeration-prime-cf-immunity}. The immunity proof
  uses only \textup{(U1)}, \textup{(U2)}, and \textup{(U3)} and does not
  assume that the full greedy language is regular. In the Pisot subclass,
  Frougny's finite
  normalization theorem supplies a finite transducer, and the associated
  canonical syntax and Bertrand/substitution descriptions are standard
  \cite{Frougny1992,BruyereHansel1997,Fabre1995,FrougnySolomyak1992}.
  Finally, multiple context-free, context-free, and
  regular languages are closed under reversal, so changing between the LSD-first proof
  convention and the usual MSD-first convention preserves every
  conclusion.
\end{proof}

\begin{example}[Pisot and non-Pisot Perron systems]
  \label{ex:pisot-unit-nonunit}
  Fibonacci, Pell, and Tribonacci numeration have unit trailing
  coefficient, so the unit-system conclusions apply. In ordinary base
  \(b\), \(U_n=b^n\) and \(a_0=b\); the prime language is MCF-immune,
  while the unit-system iteration requires \(\lvert a_0\rvert=1\).

  There is also an infinite family of genuinely nonintegral weak-Perron
  systems. Let \(p,q\ge2\), put \(B=pq\), and suppose that \(B\) is not a
  square. Define
  \[
    U_{2t}=B^t,\qquad U_{2t+1}=pB^t\qquad(t\ge0).
  \]
  The successive quotients are alternately \(p\) and \(q\). Repeated
  Euclidean division in these alternating radices gives every positive
  integer a unique greedy expansion: the digit at an even position lies
  in \(\{0,\ldots,p-1\}\), and the digit at an odd position lies in
  \(\{0,\ldots,q-1\}\). Hence the LSD-first canonical language is regular,
  since a finite automaton need only track position parity, the relevant
  digit bound, and a nonzero final digit. Moreover,
  \[
    U_{n+2}=BU_n,
  \]
  and the irreducible polynomial \(X^2-B\) is the minimal recurrence
  polynomial. Its positive root \(\sqrt B\) is nonintegral and weak Perron,
  with conjugate \(-\sqrt B\) of the same modulus. The regular ray
  \(\{0^{2t}1:t\ge0\}\) has values \(B^t\). For example, \(p=2,q=3\)
  gives \(1,2,6,12,36,72,\ldots\).

  The polynomial \(X^2-5X+5\) gives a genuinely non-Pisot instance. Its
  dominant root \(\beta=(5+\sqrt5)/2\) is Perron, while its other conjugate
  \((5-\sqrt5)/2\) is greater than one. The integer sequence
  \[
    U_0=1,\qquad U_1=4,\qquad
    U_{n+2}=5U_{n+1}-5U_n
  \]
  begins \(1,4,15,55,200,\ldots\). Induction in the recurrence gives
  \(3<U_{n+1}/U_n\le4\), so the values increase and their quotients are
  bounded. This is a non-Pisot Perron example to which the standard greedy
  classification applies.
\end{example}

\begin{theorem}[Weak-Perron radical classification]
  \label{cor:linear-perron-bounded-support-classification}
  Let \(U=(U_n)_{n\ge0}\) be a standard greedy linear numeration basis with
  \(U_0=1\), strictly increasing integer place values with bounded
  quotients, and let \(\mathcal R_U\) be the LSD-first reversal of its
  greedy representation language. Suppose that the minimal tail recurrence
  polynomial \(\mu_U\) is the minimal polynomial of a weak Perron number
  \(\beta>1\), meaning that every conjugate \(\beta'\) satisfies
  \(\abs{\beta'}\le\beta\). Let \(S_{\mathrm r}\) be the set of prime
  divisors of \(\mu_U(0)\). The residue-growth step below is a
  coefficient-level specialization of the quotient-limit mechanism in
  \cite[Proposition~10 and Remark~12]{CharlierKreczman2025}; the
  prime-support/MCFL equivalence is not contained there. With this
  attribution understood, the following are equivalent:
  \begin{enumerate}[label=\textup{(\roman*)}]
    \item There exist \(k<\infty\) and an infinite \(k\)-MCFL
    \(L\subseteq\mathcal R_U\) such that
    \[
      \sup_{w\in L}\omega_{S_{\mathrm r}^c}(\val_U(w))<\infty;
    \]
    \item \(\mathcal R_U\) contains a geometric synchronized scheme whose
    values are \(cb^t\), with \(c\ge1\), \(b\ge2\);
    \item Some positive power of \(\beta\) is an integer:
    \[
      \beta^m=B\in\mathbb Z_{\ge2}
      \quad\text{for some }m\ge1;
    \]
    \item For some \(m\ge1\) and \(B\ge2\),
    \[
      \mu_U(X)\mid X^m-B;
    \]
    \item For some \(m\ge1\), \(B\ge2\), and \(n_0\),
    \[
      U_{n+m}=BU_n\qquad(n\ge n_0).
    \]
  \end{enumerate}
  In the positive case, \(\mathcal R_U\) already contains the infinite
  regular ray
  \[
    \{0^{n_0+mt}1:t\ge0\},\qquad
    \val_U(0^{n_0+mt}1)=U_{n_0}B^t.
  \]
  Consequently, if no positive power of \(\beta\) is an integer, then every
  infinite MCFL \(L\subseteq\mathcal R_U\) has
  \[
    \sup_{w\in L}\omega_{S_{\mathrm r}^c}(\val_U(w))=\infty.
  \]
\end{theorem}

The proof has four stages. We first identify the peripheral period of the
weak-Perron spectrum, then determine the tail coefficients and their Galois
covariance. Residue asymptotics propagate positivity through the peripheral
cycle and give the global root limit. The greedy length interval then
converts a geometric synchronized ray into an integral power of \(\beta\).

\begin{proof}
  The equivalence of \textup{(i)} and \textup{(ii)} is
  \Cref{thm:bounded-outside-support-geometric-ray}.

    We first establish the asymptotic needed for
  \textup{(ii)}\(\Rightarrow\)\textup{(iii)}. This conclusion is also
  obtainable from Charlier--Kreczman's quotient-limit mechanism
  \cite[Proposition~10 and Remark~12]{CharlierKreczman2025}; we give the
  coefficient-level argument because the residue thresholds and positivity
  are part of the hypotheses used here.

  \emph{Peripheral period.}
  The classical power characterization of weak Perron numbers gives an
  integer \(h\ge1\) for which
  \[
    \rho:=\beta^h
  \]
  is a Perron number \cite{Brunotte2013,Lind1984}. Equivalently, if
  \(\gamma\) is any conjugate of \(\beta\), then
  \[
    \lvert\gamma\rvert=\beta\ \Longrightarrow\ \gamma^h=\rho,
    \qquad
    \lvert\gamma\rvert<\beta\ \Longrightarrow\
      \lvert\gamma^h\rvert<\rho.
  \]
  In particular, every peripheral phase \(\gamma/\beta\) is an
  \(h\)-th root of unity; the conclusion is stronger than merely placing
  \(\gamma^h\) somewhere on the circle of radius \(\rho\).

  \emph{Tail coefficients and their covariance.}
  Put \(d=\deg\mu_U\), let \(K\) be a splitting field of \(\mu_U\), and
  let \(J_{\mathrm r}\) be the threshold from which \(\mu_U\) annihilates
  \(U\). The irreducible polynomial \(\mu_U\in\mathbb Q[X]\) is
  separable. Hence there are unique coefficients \(A_\gamma\in K\), one
  for each conjugate \(\gamma\) of \(\beta\), such that
  \begin{equation}
    U_n=\sum_\gamma A_\gamma\gamma^n
    \qquad(n\ge J_{\mathrm r}).
    \label{eq:weak-perron-tail-expansion}
  \end{equation}
  Indeed, the coefficients are determined by the \(d\) values with
  \(n=J_{\mathrm r},\ldots,J_{\mathrm r}+d-1\); the corresponding matrix
  is a nonsingular Vandermonde matrix multiplied by the nonzero diagonal
  matrix \(\operatorname{diag}(\gamma^{J_{\mathrm r}})\). No polynomial
  factors in \(n\) occur because every root is simple.

  For \(\sigma\in\operatorname{Gal}(K/\mathbb Q)\), applying \(\sigma\)
  to \eqref{eq:weak-perron-tail-expansion} and using
  \(U_n\in\mathbb Q\) gives
  \[
    U_n=\sum_\gamma \sigma(A_\gamma)\sigma(\gamma)^n.
  \]
  Uniqueness of the expansion therefore yields the explicit covariance
  identity
  \begin{equation}
    A_{\sigma(\gamma)}=\sigma(A_\gamma).
    \label{eq:weak-perron-galois-covariance}
  \end{equation}
  For each conjugate \(\gamma\), the embedding
  \(\mathbb Q(\beta)\hookrightarrow K\) that sends \(\beta\) to
  \(\gamma\) extends to an automorphism of the splitting field \(K\).
  Thus, if \(A_\beta=0\),
  \eqref{eq:weak-perron-galois-covariance} forces every \(A_\gamma\) to vanish,
  contradicting the positive, nonzero tail of \(U\). Consequently
  \(A_\beta\ne0\); the same argument in fact shows that every
  \(A_\gamma\ne0\).

  \emph{Residue asymptotics.}
  For each \(r\in\{0,\ldots,h-1\}\), set
  \begin{equation}
    N_r:=\max\!\left\{0,
      \left\lceil\frac{J_{\mathrm r}-r}{h}\right\rceil\right\},
    \qquad
    C_r:=\sum_{\lvert\gamma\rvert=\beta}A_\gamma\gamma^r.
    \label{eq:weak-perron-residue-threshold-coefficient}
  \end{equation}
  Then \(hn+r\ge J_{\mathrm r}\) whenever \(n\ge N_r\), so
  \eqref{eq:weak-perron-tail-expansion} applies. For a peripheral conjugate,
  \(\gamma^{hn+r}=\gamma^r\rho^n\).
  If \(\lvert\gamma\rvert<\beta\), then
  \(\lvert\gamma^h\rvert<\rho\). If such conjugates exist, their total
  contribution is \(O(\theta^n)\), where
  \(\theta=\max_{\lvert\gamma\rvert<\beta}\lvert\gamma\rvert^h<\rho\);
  if none exist, that contribution is identically zero and we choose any
  \(0<\theta<\rho\). Therefore, for
  every fixed residue \(r\),
  \begin{equation}
    U_{hn+r}=C_r\rho^n+O(\theta^n)
             =C_r\rho^n+o(\rho^n)
    \qquad(n\ge N_r).
    \label{eq:weak-perron-residue-asymptotic}
  \end{equation}
  In particular,
  \[
    C_r=\lim_{n\to\infty}\rho^{-n}U_{hn+r}.
  \]
  The right-hand side is a limit of real nonnegative numbers, so
  \(C_r\in\mathbb R\) and \(C_r\ge0\).

  At least one residue coefficient is nonzero. To see this without any
  suppressed independence argument, list the peripheral conjugates as
  \(\gamma_1,\ldots,\gamma_s\), and write
  \(\gamma_j=\beta\zeta_j\). The \(\zeta_j\) are distinct \(h\)-th
  roots of unity, so \(s\le h\). If \(C_r=0\) for every
  \(0\le r<h\), then
  \begin{equation}
    0=\beta^{-r}C_r
      =\sum_{j=1}^s A_{\gamma_j}\zeta_j^r
      \qquad(0\le r<h).
      \label{eq:weak-perron-fourier-system}
  \end{equation}
  The \(h\times s\) Fourier--Vandermonde matrix
  \[
    (\zeta_j^r)_{\substack{0\le r<h\\1\le j\le s}}
  \]
  has full column rank: its first \(s\) rows form a square Vandermonde
  matrix with determinant
  \(\prod_{1\le i<j\le s}(\zeta_j-\zeta_i)\ne0\).
  Equation \eqref{eq:weak-perron-fourier-system} would force all
  peripheral coefficients to vanish,
  including \(A_\beta\), a contradiction.

  \emph{Cyclic propagation of positivity.}
  Choose \(r\) with \(C_r>0\). If \(r<h-1\), then for all sufficiently
  large \(n\), strict increase gives
  \[
    U_{hn+r}<U_{hn+r+1}.
  \]
  Divide by \(\rho^n\) and use
  \eqref{eq:weak-perron-residue-asymptotic} for both residue classes. The
  limiting inequality gives \(C_r\le C_{r+1}\), so in particular
  \(C_{r+1}>0\). At the wraparound, strict increase gives
  \begin{equation}
    U_{hn+h-1}<U_{h(n+1)},
    \qquad
    \rho^{-n}U_{h(n+1)}\longrightarrow\rho C_0,
    \label{eq:weak-perron-wraparound-limit}
  \end{equation}
  whereas \(\rho^{-n}U_{hn+h-1}\to C_{h-1}\). Thus
  \(C_{h-1}\le\rho C_0\), and \(C_{h-1}>0\) forces \(C_0>0\).
  Starting from one positive coefficient and moving cyclically proves
  \begin{equation}
    C_r>0\qquad(0\le r<h).
    \label{eq:weak-perron-positive-residues}
  \end{equation}

  Finally, if \(N=hn+r\), then
  \eqref{eq:weak-perron-residue-asymptotic} and
  \eqref{eq:weak-perron-positive-residues} give
  \begin{align*}
    U_N&=C_r\rho^n(1+o(1)),\\
    \frac{\log U_N}{N}
      &=\frac{n\log\rho+\log C_r+o(1)}{hn+r}
        \longrightarrow\frac{\log\rho}{h}=\log\beta.
  \end{align*}
  There are only finitely many residues, so this proves the global limit
  \begin{equation}
    \lim_{N\to\infty}U_N^{1/N}=\beta.
    \label{eq:weak-perron-global-root-limit}
  \end{equation}


  Suppose now that \textup{(ii)} holds, and write
  \[
    \val_U(W(t))=cb^t,
    \qquad \abs{W(t)}=L+Dt,
  \]
  where \(D>0\) is the total pumped length.  In a greedy \(U\)-system, a
  canonical word of length \(m\) represents a value in the half-open interval
  \([U_{m-1},U_m)\).  Hence
  \[
    U_{L+Dt-1}\le cb^t<U_{L+Dt}.
  \]
  For \(k_t=L+Dt\), the root limit gives
  \[
    U_{k_t}^{1/t}
      =\left(U_{k_t}^{1/k_t}\right)^{k_t/t}\longrightarrow\beta^D,
    \qquad
    U_{k_t-1}^{1/t}
      =\left(U_{k_t-1}^{1/(k_t-1)}\right)^{(k_t-1)/t}
       \longrightarrow\beta^D.
  \]
  Since \((cb^t)^{1/t}\to b\), taking \(t\)-th roots in the preceding
  squeeze yields \(b=\beta^D\). This is \textup{(iii)} with \(m=D\) and
  \(B=b\).

  Conditions \textup{(iii)} and \textup{(iv)} are equivalent because
  \(\beta^m=B\) says exactly that the minimal polynomial \(\mu_U\) of
  \(\beta\) divides \(X^m-B\) in \(\mathbb Q[X]\). Conditions
  \textup{(iv)} and \textup{(v)} are equivalent by the definition of
  \(\mu_U\) as the monic generator of \(I_U^{\mathrm{tail}}\): the
  divisibility \(\mu_U\mid X^m-B\) puts \(X^m-B\) in that ideal and gives
  \(U_{n+m}=BU_n\) eventually, while an eventual identity of this form
  puts \(X^m-B\) in the ideal and hence forces the divisibility.

  Finally, assume \textup{(v)} and enlarge \(n_0\), if necessary, so that
  \(n_0\ge J_{\mathrm r}\). The greedy representation of \(U_n\), in the
  LSD-first convention, is \(0^n1\). Taking
  \[
    a=0^{J_{\mathrm r}},\quad k=1,\quad
    u_1=0^{n_0-J_{\mathrm r}},\quad v_1=0^m,\quad
    w_1=s_1=\varepsilon,\quad u_2=1
  \]
  gives canonical words \(W(t)=0^{n_0+mt}1\) and
  \[
    \val_U(W(t))=U_{n_0}B^t.
  \]
  This is the asserted regular geometric ray and proves
  \textup{(v)}\(\Rightarrow\)\textup{(ii)}. The final assertion is the
  negation of \textup{(i)}. In the strict-Perron subclass, the equality
  \(b=\beta^D\) forces every conjugate to have modulus \(\beta\) unless
  \(\beta\) has degree one; hence the theorem recovers the previous
  integral-root boundary there.
\end{proof}

\begin{remark}[Scope]
  \label{rem:pisot-open-interfaces}
  The results concern integer-valued greedy linear numeration with an
  eventual integral recurrence. Real-base dynamical expansions and
  genealogical-rank abstract numeration use different evaluation maps. For
  a nonunit recurrence witness, \Cref{thm:nonunit-mcfl-escape-dichotomy}
  retains the alternative between additional primes and increasing valuations;
  within the weak-Perron greedy subclass,
  \Cref{cor:linear-perron-bounded-support-classification} identifies the
  latter bounded-support regime by integrality of a positive power of the
  dominant root.
\end{remark}
